\section{从回答到行动：Agent OS 与对齐要求}

本章把自动驾驶问题推广到 Agent。李想认为，五年内未必会有通用 Agent，但会有 Agent OS。这个判断的含义是：不同专业工种需要不同 Agent，但这些 Agent 需要共同运行在一个操作系统式底座里，获得工具、权限、上下文、虚拟机、记忆和安全边界。Agent OS 不是一个聊天窗口，而是专业 Agent 的基础设施。

\begin{figure}[H]
\centering
\includegraphics[width=0.96\textwidth]{action-agent-alignment.png}
\caption{从回答到行动：AI 进入资产、工具和车辆时，对齐要求会变高。自制概念图，依据 00:53:58--00:55:00 与 00:55:00--00:55:25 对谈内容整理。}
\end{figure}

\begin{knowledgebox}{读图：行动会放大风险}
信息检索错了，可以重搜；研究错了，可以复核；但如果 Agent 调用资产、工具、电脑、车辆或财产，错误会直接改变世界。因此 Agent 的对齐要求、权限边界和可追责性都比聊天模型更高。
\end{knowledgebox}

\subsection{为什么通用 Agent 不如专业 Agent}

李想认为，真实生产工具要替代和解放人类的高频工作，例如开车、编程、运营、销售或制造。不同工种有不同数据、工具、权限、错误代价和评价标准，因此短期更可能出现专业 Agent，而不是一个全能 Agent。Agent OS 的价值在于，为这些专业 Agent 提供共同运行环境。

\begin{importantbox}{Agent OS 的定义}
Agent OS 是让专业 Agent 运行的基础系统。它至少包括上下文管理、工具调用、权限控制、虚拟机或执行环境、记忆、日志、评测和安全策略。它的价值不在于“什么都能做”，而在于让不同专业 Agent 能可靠地做事。
\end{importantbox}

\subsection{对齐从语言问题变成行动问题}

当 AI 只是回答问题时，对齐主要是内容安全、事实准确和价值偏好；当 AI 开始调用工具和资产，对齐变成操作风险。比如一个驾驶 Agent 的错误会影响道路安全；一个财务 Agent 的错误会影响资金；一个代码 Agent 的错误会影响系统稳定。对齐因此从“模型说得对不对”扩展到“模型做得可不可控”。

\begin{knowledgebox}{对齐需求的粗略分层}
\[
\text{risk} \approx \text{action scope} \times \text{asset value} \times \text{irreversibility}.
\]
其中，\(\text{action scope}\) 表示 Agent 能操作的范围，\(\text{asset value}\) 表示被调用资产的价值，\(\text{irreversibility}\) 表示错误是否容易撤销。越接近车、钱、账号、生产系统和物理设备，对齐要求越高。
\end{knowledgebox}

\subsection{本章小结}

Agent OS 是这期从自动驾驶推到通用 AI 产品的关键桥梁。它说明 AI 的下一步不是更会聊天，而是更会在有权限、有工具、有责任的环境里行动。

